\paragraph{Micro--macro tension in generative mobility modeling.}
Mobility Digital Twins require \emph{simulation validity}: sampled trajectories should cover plausible futures while preserving physically realistic macroscopic motion statistics. Our results show that these objectives are not aligned if one optimizes only point-wise error. The deterministic prior (L2 regression) achieves the lowest mean errors on the diagnostic set, but it produces a single trajectory per condition and therefore cannot represent multi-modality (Table~\ref{tab:phaseb_quick}). In contrast, diffusion-based generators improve best-of-$K$ coverage, but exhibit \emph{macroscopic shrinkage}: net displacement and speed are under-estimated, leading to lower MSD/Rog (Figure~\ref{fig:phaseb_quant}).

\paragraph{Effect of navigation field conditioning under a leakage-free protocol.}
Navigation field conditioning (train-only) improves the oracle upper bound (best-of-$K$) relative to data-only diffusion (Table~\ref{tab:phaseb_quick}), indicating that local mean-flow information helps route-level coverage. However, it does not eliminate macroscopic shrinkage. This implies that, under our leakage-free setting, adding a mean-flow conditioning signal alone is insufficient to restore the net displacement and speed scale required for simulation validity. This observation aligns with the intuition that field-based guidance can introduce an implicit bias toward the mean (mean-reversion), dampening the stochastic expansion needed to match true macroscopic dispersion.

\paragraph{Residual diffusion as a structural fix.}
Residual decomposition separates responsibilities: the deterministic prior carries OD-driven macroscopic motion scale, while diffusion focuses on stochastic residual deviations. This directly targets the shrinkage definition used in this paper (under-estimated net displacement and speed). Empirically, residual diffusion restores macroscopic scale (Speed Ratio and Rog Ratio close to 1) while maintaining strong best-of-$K$ behavior (Table~\ref{tab:residual_test}).

\paragraph{CFG as a protocolized inference knob.}
Even after the structural fix, practical digital-twin usage requires \emph{controllability}: planners may prefer higher micro precision for typical scenarios or stronger macro-validity for stress testing and scenario analysis. CFG provides a simple inference-time knob to trade micro and macro objectives without retraining. Crucially, we avoid a parameter swamp by pre-registering two settings (cfg=2 for micro-optimal within a macro gate, cfg=3 for macro-validity emphasis), and we report the full trend only as a transparent validation curve (Figure~\ref{fig:cfg_pareto}).

\paragraph{Limitations and future work.}
Our current pipeline is grid-based. For geographic presentation we use a linear bbox projection to latitude/longitude and overlay Shenzhen district boundaries from GeoJSON; this improves interpretability but is not road-level map matching. Future work should incorporate road-graph structure and explore more expressive conditioning mechanisms (e.g., ControlNet-style injection or learnable gating) so that navigation field information guides direction without acting as an overly restrictive constraint. Finally, stronger conclusions require multi-seed replication and broader evaluation under the same strict data contract.
